\section{Auxiliary scope evidence and discussion}

Auxiliary scope evidence is limited but not empty. On Llama-3.1-8B, the MHST effect remained same-sign on held-out test with mean $\DeltaBpair = 0.9231$ and 95\% CI $[0.5897, 1.2827]$, although coverage was much lower than in the primary Qwen run. RC-ICL was also same-sign on its main barrier summaries ($0.5244$, 95\% CI $[0.2043, 0.8401]$) and kept $k{=}0$ exactly null, but it failed its no-conflict cleanliness threshold badly, reaching only $0.70$ where the locked criterion required $0.95$. The safe conclusion is therefore conservative: the effect is not obviously unique to one exact MHST/model configuration, but we do not claim broad generalization beyond the main MHST/Qwen setting. Appendix~\ref{app:scope} gives the corresponding auxiliary curves and summary table.

The paper's main contribution is behavioral and methodological. The matched-history revision-barrier construction gives a concrete way to separate revisability from current designated support state. In the primary MHST/Qwen setting, it yields a strong, well-controlled asymmetry: committed histories are harder to reverse than fresh histories even when the designated incumbent/challenger state and evidence margin are matched by construction. The state-only insufficiency decomposition sharpens that point further: even a simple monotone predictor built from designated margin and prefix-only output margin leaves a positive residual paired history effect.

The internal sections then mark the current interpretive limits. Prefix-boundary state contains signal related to later revisability, but not enough to justify a stronger latent-state claim over output margin. One-site steering does not provide a usable causal handle under the locked protocol, including after a principled boundary-alignment rerun. These follow-ups therefore increase trust not by rescuing a stronger mechanism story, but by showing exactly what the current evidence does \emph{not} support.

The main limitations are straightforward. MHST is synthetic and deliberately controlled, so the strongest claim is about revisability under matched designated state rather than about broad real-world generalization. Auxiliary scope evidence remains limited. And the negative causal result closes only the current one-site steering operationalization; it does not show that no stronger causal account could exist. With those caveats in place, the final claim is narrow but substantive: the matched-history revision-barrier construction is a useful behavioral diagnostic for history-sensitive revisability under controlled designated state.
